%%%%%%%%%%  scan idx 279  |  folio 266  %%%%%%%%%%
\SGAPrintPage{266}%
If \(E\) satisfies condition (5.1 a)) for the integer \(p\), \(E^{\vee}\) satisfies the same condition for the integer \(-p-1\) by (5.2 (iii)), and the arrow \((*)\) is an isomorphism by the first assertion of (5.3.1). This completes the proof of (5.3.1), and hence of (5.3).

\sgaStableTarget{sga6.e03.rem.5-3-2}\textbf{Remark 5.3.2.} Proposition (5.3) applies in particular when \(E\) is a complex of bounded degrees of quasi-coherent \(\underline{O}_X\)-Modules that are flat and of finite presentation rel. to \(f\): indeed, \(E\) then has finite perfect amplitude rel. to \(f\) by (4.3.2).

%%%%%%%%%%  scan idx 280  |  folio 267  %%%%%%%%%%
\SGAPrintPage{267}%
\sgaStableTarget{sga6.e03.def.5-4}\textbf{Definition 5.4.} \uline{Let} \(X\) \uline{be a topos}, \uline{and denote by} \(\pt(X)\) \uline{the set of its points} (SGA 4 IV 6.1). \uline{Let} \(M\) \uline{be an ordered set}, \(f:\pt(X)\to M\) \uline{a map}, \uline{and} \(x\in\pt(X)\). \uline{We say that} \(f\) \uline{is upper semi-continuous} (\uline{resp.} \uline{lower semi-continuous}, \uline{resp.} \uline{locally constant}) \uline{at} \(x\) \uline{if there exists} \uline{an} \(x\)-\uline{pointed object} \(U\) \uline{of} \(X\) \uline{such that for every point} \(y\) \uline{of} \(U\) \uline{we have} \(f(y)\leq f(x)\) (\uline{resp.} \(f(y)\geq f(x)\), \uline{resp.} \(f(y)=f(x)\)). \uline{We say that} \(f\) \uline{is upper semi-continuous} (\uline{resp.} \uline{lower semi-continuous}, \uline{resp.} \uline{locally constant}) \uline{if it is so at every point}\SGAAuthorFootnote{1}{When \(X\) comes from a topological space, we recover the usual notions.}.

\sgaStableTarget{sga6.e03.ex.5-4-1}\textbf{Examples 5.4.1.} a) Denote by \(M_X\) the constant sheaf with value \(M\) on \(X\), and let \(s\in H^0(X,M)\). The map \(x\mapsto s(x)=x^*(s)\) from \(\pt(X)\) to \(M\) is locally constant. In particular, let \(E\) be a perfect complex on \(X\), supposed locally ringed: the rank of \(E\) (I 6.11) defines a locally constant function on \(\pt(X)\) with values in \(\mathbb Z\).

b) Suppose \(X\) locally ringed, and let \(P\) be an \(\underline{O}_X\)-Module of finite type. For \(x\in\pt(X)\), denote by \(\underline{O}_{X,x}\) the local stalk ring of \(\underline{O}_X\) at \(x\), and by \(k(x)\) its residue field. Then, for \(x\in\pt(X)\), \(P_x\otimes_{\underline{O}_{X,x}}k(x)\) has finite rank over \(k(x)\), and the function \(x\mapsto\operatorname{rg}_{k(x)}(P\otimes_{\underline{O}_{X,x}}k(x))\) is upper semi-continuous. Indeed, if, for \(x\in\pt(X)\), \(P_x\otimes_{\underline{O}_{X,x}}k(x)\) has rank \(n\), there is an epimorphism \(\underline{O}_{X,x}^n\to P_x\) by Nakayama, and, since \(P\) is of finite type, there is an \(x\)-pointed object \(U\) of \(X\) and an epimorphism \(\underline{O}_U^n\to P\,|\,U\), giving the assertion.

\sgaStableTarget{sga6.e03.lem.5-5}\textbf{Lemma 5.5.} \uline{Let} \((X,\underline{O}_X)\) \uline{be a locally ringed topos}, \(E\in\Ob D^{-}(X)\), \uline{and} \(p\in\mathbb Z\). \uline{Then we have}, \uline{with the notation of} (5.4.1 b)):
\begin{enumerate}[label=(\roman*)]
%%%%%%%%%%  scan idx 281  |  folio 268  %%%%%%%%%%
\item\SGAPrintPage{268}%
\uline{If} \(E\) \uline{is perfect}, \uline{the} \(H^i(E_x\otimes_{\underline{O}_{X,x}}^{\mathrm L}k(x))\), \uline{for} \(x\in\pt(X)\), \uline{are vector spaces of finite rank over} \(k(x)\), \uline{all but finitely many zero}, \uline{and the function}
\[
 x\longmapsto\sum(-1)^i\operatorname{rg}_{k(x)}H^i(E_x\otimes_{\underline{O}_{X,x}}^{\mathrm L}k(x))
\]
\uline{is locally constant on} \(\pt(X)\).
\item \uline{If} \(E\) \uline{is} \((p-1)\)-\uline{pseudo-coherent}, \(H^p(E_x\otimes_{\underline{O}_{X,x}}^{\mathrm L}k(x))\) \uline{is}, \uline{for} \(x\in\pt(X)\), \uline{a vector space of finite rank over} \(k(x)\), \uline{and the function}
\[
 x\longmapsto\operatorname{rg}_{k(x)}H^p(E_x\otimes_{\underline{O}_{X,x}}^{\mathrm L}k(x))
\]
\uline{is upper semi-continuous on} \(\pt(X)\). \uline{If moreover the functor} \(M\mapsto H^p(E\otimes^{\mathrm L}M)\) \uline{from} \(\Mod(X)\) \uline{to} \(\Mod(X)\) \uline{is exact}, \(H^p(E)\) \uline{is locally free of finite type and the preceding function is locally constant}.
\item \uline{Suppose} \(E\) \(p\)-\uline{pseudo-coherent} (\uline{resp.} \((p-1)\)-\uline{pseudo-coherent}), \uline{and let} \(x\in\pt(X)\). \uline{If the functor} \(M\mapsto H^p(E_x\otimes_{\underline{O}_{X,x}}^{\mathrm L}M)\) \uline{is right exact} (\uline{resp.} \uline{left exact}) \uline{on the category} \(\Mod(x)\) \uline{of} \(\underline{O}_{X,x}\)-\uline{Modules}, \uline{there is an} \(x\)-\uline{pointed open} \(U\) \uline{of} \(X\) \uline{such that the functor} \(M\mapsto H^p(E\,|\,U\otimes^{\mathrm L}M)\) \uline{is right exact} (\uline{resp.} \uline{left exact}) \uline{on} \(\Mod(U)\).
\end{enumerate}

\uline{Proof}. (i) By (I 4.19.1), \(E_x\otimes_{\underline{O}_{X,x}}^{\mathrm L}k(x)\) is a perfect complex of \(k(x)\)-modules, so (I 5.9 and 3.5) the \(H^i(E_x\otimes_{\underline{O}_{X,x}}^{\mathrm L}k(x))\) have finite rank over \(k(x)\) and all but finitely many are zero. We have
\begin{equation}
 \operatorname{rg}(E_x\otimes_{\underline{O}_{X,x}}^{\mathrm L}k(x))=\sum(-1)^i\operatorname{rg}_{k(x)}H^i(E_x\otimes_{\underline{O}_{X,x}}^{\mathrm L}k(x)),\qquad\mbox{\SGAAuthorFootMark{1}}\tag{5.5.1}
\end{equation}
\SGAAuthorFootText{1}{\(\operatorname{rg}(\ )\) denotes the rank of a perfect complex (I 6.11).}%
as is seen by an immediate dévissage using the additivity of the rank function. In view of (5.4.1 a)), assertion (i) therefore follows from
%%%%%%%%%%  scan idx 282  |  folio 269  %%%%%%%%%%
\SGAPrintPage{269}%
the compatibility of rank with inverse images (I 6.12).

(ii) After truncating and shifting, we may, for the first assertion, suppose \(E\) strictly perfect, \(p=0\), and \(E^i=0\) for \(i\notin[-1,+1]\). For \(x\in\pt(X)\), put \(E_x\otimes_{\underline{O}_{X,x}}^{\mathrm L}k(x)=E_{k(x)}\). By (5.5.1) we have
\[
\begin{gathered}
 \operatorname{rg}_{k(x)}H^0(E_{k(x)})=\operatorname{rg}_{k(x)}H^{-1}(E_{k(x)})+\operatorname{rg}_{k(x)}H^1(E_{k(x)})\\
 +\operatorname{rg}(E_x),
\end{gathered}
\]
so it suffices to prove that the function \(x\mapsto\operatorname{rg}_{k(x)}H^{-1}(E_{k(x)})\) (resp. \(\operatorname{rg}_{k(x)}H^1(E_{k(x)})\)) is upper semi-continuous (since the function \(x\mapsto\operatorname{rg}(E_x)=\operatorname{rg}(E_{k(x)})\) is locally constant). For the function \(x\mapsto\operatorname{rg}_{k(x)}H^1(E_{k(x)})\), this is an immediate consequence of (5.4.1 b))\SGAAuthorFootnote{*}{taking into account that \(H^1(E_{k(x)})\simeq H^1(E)_{k(x)}\).}. For the function \(x\mapsto\operatorname{rg}_{k(x)}H^{-1}(E_{k(x)})\), we may suppose \(E^1=0\), i.e. replace \(E\) by \(E'=(0\to E^{-1}\to E^0\to0)\). We already know that the function \(x\mapsto\operatorname{rg}_{k(x)}H^0(E'_{k(x)})\) is upper semi-continuous, so the same holds for the function \(x\mapsto\operatorname{rg}_{k(x)}H^{-1}(E'_{k(x)})=\operatorname{rg}_{k(x)}H^{-1}(E_{k(x)})\), since the function \(x\mapsto\operatorname{rg}_{k(x)}H^0(E'_{k(x)})-\operatorname{rg}_{k(x)}H^{-1}(E'_{k(x)})=\operatorname{rg}(E'_x)\) is locally constant. Let us now prove the second assertion. After truncating, we may suppose \(E\) strictly perfect and \(E^i=0\) for \(i<p-1\). The \(\underline{O}_X\)-Module \(Z'^i(E)=\Coker(E^{i-1}\to E^i)\) is therefore of finite presentation. We have an exact sequence:
\begin{equation}
 0\to H^p(E)\to Z'^p(E)\to E^{p+1}\to Z'^{p+1}(E)\to0.\tag{$*$}
\end{equation}
Since the functor \(M\mapsto H^p(E\otimes^{\mathrm L}M)\) is right and left exact, \(Z'^{p+1}(E)\) and \(Z'^p(E)\) are flat by (5.2 (iii)), hence (I 5.8.3) locally free of finite type. The exactness of (*) then implies that \(H^p(E)\) is locally free of finite type, and consequently that the function \(x\mapsto\operatorname{rg}_{k(x)}H^p(E_{k(x)})=\operatorname{rg}(H^p(E)_x)\)
%%%%%%%%%%  scan idx 283  |  folio 270  %%%%%%%%%%
\SGAPrintPage{270}%
is locally constant, which completes the proof of (ii).

(iii) We may suppose \(E\) strictly perfect and \(E^i=0\) for \(i<p\) (resp. \(p-1\)). Saying that the functor \(M\mapsto H^p(E_x\otimes_{\underline{O}_{X,x}}^{\mathrm L}M)\) is right exact (resp. left exact) on \(\Mod(x)\) is equivalent, by (5.2 (iii)), to saying that \(Z'^{p+1}(E_x)\) (resp. \(Z'^p(E_x)\)) (with the notation of (ii)) is flat, hence free of finite type (since it is of finite presentation). Since \(Z'^{p+1}(E_x)\simeq Z'^{p+1}(E)_x\) (resp. \(Z'^p(E_x)\simeq Z'^p(E)_x\)), the usual argument (cf. EGA \(\mathrm O_I\) 5.2.7) shows that there is an \(x\)-pointed object \(U\) of \(X\) such that \(Z'^{p+1}(E\,|\,U)\) (resp. \(Z'^p(E\,|\,U)\)) is free of finite type, and the result follows.

\sgaStableTarget{sga6.e03.cor.5-6}\textbf{Corollary 5.6} (EGA III 7.6.9). \uline{Let} \(X\) \uline{be a scheme}, \(p\in\mathbb Z\), and \(E\in\Ob D^{-}(X)\) \uline{a} \((p-1)\)-\uline{pseudo-coherent complex}. \uline{Then}:
\begin{enumerate}[label=(\roman*)]
\item \uline{For} \(x\in X\), \(H^p(E\otimes_X^{\mathrm L}k(x))\) \uline{has finite rank over} \(k(x)\) \uline{and the function}
\[
 x\longmapsto\operatorname{rg}_{k(x)}H^p(E\otimes_X^{\mathrm L}k(x))
\]
\uline{is constructible and upper semi-continuous on} \(X\).
\item \uline{Suppose} \(X\) \uline{quasi-separated and} \(E\) \uline{with quasi-coherent and locally bounded cohomology}. \uline{If the functor} \(M\mapsto H^p(E\otimes^{\mathrm L}M)\) \uline{is exact on the category of quasi-coherent} \(\underline{O}_X\)-\uline{Modules}, \uline{the function} \(x\mapsto\operatorname{rg}_{k(x)}H^p(E\otimes_X^{\mathrm L}k(x))\) \uline{is locally constant and} \(H^p(E)\) \uline{is locally free of finite type}. \uline{Conversely}, \uline{if} \(X\) \uline{is locally noetherian and reduced and the preceding function is locally constant}, \uline{the functor} \(M\mapsto H^p(E\otimes^{\mathrm L}M)\) \uline{is exact on the category of quasi-coherent} \(\underline{O}_X\)-\uline{Modules}.
\end{enumerate}

\uline{Proof}. If \(X\) is integral with generic point \(x\), the functor \(M\mapsto H^p(E_x\otimes_{\underline{O}_{X,x}}^{\mathrm L}M)\) is exact on the category of \(\underline{O}_{X,x}\)-modules (since \(\underline{O}_{X,x}=k(x)\)), so (5.5 (iii)) there is an open neighborhood \(U\) of \(x\) such that the functor \(M\mapsto H^p(E\,|\,U\otimes^{\mathrm L}M)\) is exact on \(\Mod(U)\), and, on \(U\), the function
